\subsection{Examples of partial operators}

In this section, K = R or C.

Examples. --- 1) Let X be a locally compact topological space and let \(\mu\) a positive measure on X. We fix elements \(p_1\) and \(p_2\) of \([1, +\infty[\) .

Let \(g\) a function \(\mu\) -measurable on \(X\) with values in \(K\) . Denote by \(D\) the subspace of \(\mathcal{L}_{K}^{p_1}(X, \mu)\) consisting of functions \(f\) in \(\mathcal{L}_{K}^{p_1}(X, \mu)\) such that \(gf \in \mathcal{L}_{K}^{p_2}(X, \mu)\) . The linear mapping from \(D\) in \(\mathcal{L}_{K}^{p_2}(X, \mu)\) defined by \(f \mapsto gf\) determines a partial operator of \(\mathcal{L}_{K}^{p_1}(X, \mu)\) in \(\mathcal{L}_{K}^{p_2}(X, \mu)\) , which we denote \(m_g\) .

The vector subspace of functions \(\mu\) -negligible in \(\mathcal{L}_{\mathrm{K}}^{p_1}(\mathrm{X},\mu)\) is contained in D, and the image under \(m_g\) of a function \(\mu\) -negligible is still \(\mu\) -negligible. We will denote \(\tilde{m}_g\) the partial operator from \(\mathrm{L}_{\mathrm{K}}^{p_1}(\mathrm{X},\mu)\) in \(\mathrm{L}_{\mathrm{K}}^{p_2}(\mathrm{X},\mu)\) deduced from \(m_g\) by passing to quotients. This is called the operator of multiplication by \(g\) of \(\mathrm{L}_{\mathrm{K}}^{p_1}(\mathrm{X},\mu)\) in \(\mathrm{L}_{\mathrm{K}}^{p_2}(\mathrm{X},\mu)\) . Of the functions \(g_1\) and \(g_2\) locally equal \(\mu\) -almost everywhere define the same multiplication operator.

PROPOSITION 5. — The multiplication operator \(\widetilde{m}_{g}\) of \(\mathrm{L}_{\mathrm{K}}^{p_{1}}(\mathrm{X},\mu)\) in \(\mathrm{L}_{\mathrm{K}}^{p_{2}}(\mathrm{X},\mu)\) is a closed operator with dense domain.

Let us first prove that the partial operator \(\widetilde{m}_g\) is closed. Let \((f_n, h_n)_{n \in \mathbf{N}}\) be a sequence in \(\mathcal{L}_{\mathrm{K}}^{p_1}(\mathrm{X}, \mu) \times \mathcal{L}_{\mathrm{K}}^{p_2}(\mathrm{X}, \mu)\) such that the sequence \((\widetilde{f}_n, \widetilde{h}_n)\) of the classes of \(f_n\) and \(h_n\) belongs to the graph of \(\widetilde{m}_g\) and converges in \(\mathrm{L}_{\mathrm{K}}^{p_1}(\mathrm{X}, \mu) \times \mathrm{L}_{\mathrm{K}}^{p_2}(\mathrm{X}, \mu)\) as \(n\) tends to infinity. Let \((f, h)\) be a pair in \(\mathcal{L}_{\mathrm{K}}^{p_1}(\mathrm{X}, \mu) \times \mathcal{L}_{\mathrm{K}}^{p_2}(\mathrm{X}, \mu)\) such that the pair \((\widetilde{f}, \widetilde{h})\) of their classes is the limit of \((\widetilde{f}_n, \widetilde{h}_n)\) .

There exists a sequence \((f_{n_{k}})_{k\in\mathbf{N}}\) extracted from the sequence \((f_{n})_{n}\) such that \(f_{n_{k}}(x)\) converges to \(f(x)\) for \(\mu\) -almost every x (INT, IV, p. 131, § 3, no. 4, th. 3). This implies that \(h_{n_{k}}(x)=g(x)f_{n_{k}}(x)\) converges \(\mu\) -almost everywhere to \(g(x)f(x)\) . Moreover, the sequence \((h_{n_{k}})\) converges to h in the space \(\mathcal{L}_{\mathrm{K}}^{p_{2}}(\mathrm{X},\mu)\) . The functions h and gf are therefore equal \(\mu\) -almost everywhere (loc. cit.). Thus \(\widetilde{f}\) belongs to the domain of \(\widetilde{m}_{g}\) and \(\widetilde{h}=\widetilde{m}_{g}(\widetilde{f})\) . This proves that \(\widetilde{m}_{g}\) is closed.

Let us prove that the domain of \(\widetilde{m}_g\) is dense in \(\mathrm{L}_{\mathrm{K}}^{p_1}(\mathrm{X},\mu)\) . It suffices to verify that the classes of functions \(f\in \mathcal{K}(\mathrm{X};\mathrm{K})\) belong to the closure of the domain of \(\widetilde{m}_g\) in \(\mathrm{L}_{\mathrm{K}}^{p_1}(\mathrm{X},\mu)\) Let \(f\in \mathcal{K}(\mathrm{X};\mathrm{K})\) and let \(\widetilde{f}\) its class in \(\mathrm{L}_{\mathrm{K}}^{p_1}(\mathrm{X},\mu)\) . For every integer \(n\in \mathbf{N}\) , denote \(\varphi_{n}\) the characteristic function of the set of elements \(x\in \mathrm{X}\) such that \(|g(x)|\leqslant n\) and we set \(f_{n} = f\varphi_{n}\) . We then have \(|gf_n|\leqslant n|f|\) , which belongs to \(\mathcal{L}_{\mathrm{K}}^{p_2}(\mathrm{X},\mu)\) , hence \(f_{n}\) belongs to the domain of \(\widetilde{m}_g\) . For every element \(x\) of X, the sequence \((f_n(x))_{n\in \mathbf{N}}\) converges to \(f(x)\) as \(n\to +\infty\) ; moreover, we have \(|f_n|\leqslant |f|\) , which belongs to \(\mathcal{L}_{\mathrm{K}}^{p_1}(\mathrm{X},\mu)\) . By Lebesgue's theorem (INT, IV, p. 137, § 3, n°7, th. 6), the sequence of classes of \(f_{n}\) converges to \(\widetilde{f}\) in \(\mathrm{L}_{\mathrm{K}}^{p_1}(\mathrm{X},\mu)\) . Thus, the class of \(f\) belongs to the closure of the domain of \(\widetilde{m}_g\) .

In the following proposition, we assume that \(p_1 = p_2 = 2\) .

PROPOSITION 6. --- a) Let \(g'\) a function \(\mu\) -measurable on X such that \(|g|\leqslant |g'|\) . We have \(\mathrm{dom}(\widetilde{m}_{g'})\subset\mathrm{dom}(\widetilde{m}_{g})\) and \(\mathrm{dom}(\widetilde{m}_{g'})\) is a core of the partial operator \(\widetilde{m}_{g}\) ;

\begin{enumerate}
\def\labelenumi{\alph{enumi})}
\setcounter{enumi}{1}
\tightlist
\item
  Let F be a subspace of \(\mathcal{L}_{\mathrm{K}}^{2}(\mathrm{X},\mu)\) whose intersection with \(\mathcal{K}(\mathrm{X};\mathrm{K})\) is dense in \(\mathcal{K}(\mathrm{X};\mathrm{K})\) and whose image G in \(\mathrm{L}_{\mathrm{K}}^{2}(\mathrm{X},\mu)\) is contained in \(\mathrm{dom}(\widetilde{m}_g)\) . If \(|g|^2\) is locally \(\mu\) -integrable, then \(\mathcal{K}(\mathrm{X};\mathrm{K})\) is contained in the domain of \(\widetilde{m}_g\) and G is a core of \(m_g\) .
\end{enumerate}

Let us prove \(a\) ). If \(f \in \mathcal{L}_{\mathrm{K}}^{2}(\mathrm{X}, \mu)\) belongs to the domain of \(m_{g'}\) , so that \(fg' \in \mathcal{L}_{\mathrm{K}}^{2}(\mathrm{X}, \mu)\) , the hypothesis implies that \(fg \in \mathcal{L}_{\mathrm{K}}^{2}(\mathrm{X}, \mu)\) , whence the result.

Let us prove that \(\mathrm{dom}(\widetilde{m}_{g'})\) is a core of \(\widetilde{m}_g\) , that is, the domain of \(\widetilde{m}_{g'}\) is dense in the Hilbert space \(\mathrm{E}_{\widetilde{m}_g}\) Let \(h \in \mathcal{L}_{\mathrm{K}}^2(\mathrm{X}, \mu)\) whose class \(\widetilde{h}\) belongs to \(\mathrm{E}_{\widetilde{m}_g}\) and is orthogonal to \(\mathrm{dom}(\widetilde{m}_{g'})\) This means that

\[
(\widetilde {h} \mid \widetilde {h} ^ {\prime}) _ {\widetilde {m} _ {g}} = \int_ {\mathrm{X}} \overline {{h}} h ^ {\prime} (1 + | g | ^ {2}) d \mu = 0
\]

for every function \(h' \in \mathcal{L}_{\mathrm{K}}^{2}(\mathrm{X}, \mu)\) whose class \(\widetilde{h}'\) belongs to \(\operatorname{dom}(\widetilde{m}_{g'})\) .

Let C be a compact subset of X and let \(\varphi\) its characteristic function. Let \(n \in \mathbf{N}\) . Let \(\varphi_n\) be the characteristic function of the set \(\mu\) -integrable \(\mathrm{C}_n\) of the \(x \in \mathrm{C}\) such that \(|h(x)| \leqslant n\) and we set \(h_n' = \varphi_n h\) . The class of \(h_n'\) belongs to the domain of \(\tilde{m}_{g'}\) since \(|g'h_n'| \leqslant n\varphi\) ; it follows that

\[
0 = \int_ {\mathrm{X}} \overline {{h}} h _ {n} ^ {\prime} (1 + | g | ^ {2}) d \mu = \int_ {\mathrm{X}} | h | ^ {2} \varphi_ {n} (1 + | g | ^ {2}) d \mu .
\]

This implies that h is zero for \(\mu\) -almost all \(x \in C_{n}\) and so, since n is arbitrary, that h is zero for \(\mu\) -almost all \(x \in C\) It finally follows that h is zero \(\mu\) \(\mu\)-almost everywhere, since C is arbitrary and h is moderate (INT, V, p. 9, § 1, n°3, cor.).

Consider now the assertion \(b\) ). Since \(|g|^2\) is locally \(\mu\) -integrable, the function \(fg\) belongs to \(\mathcal{L}_{\mathrm{K}}^{2}(\mathrm{X},\mu)\) if \(f\in\mathcal{K}(\mathrm{X};\mathrm{K})\) , hence \(\mathcal{K}(\mathrm{X};\mathrm{K})\) is contained in the domain of \(\widetilde{m}_{g}\) .

Let \(h \in \mathcal{L}_{\mathrm{K}}^{2}(\mathrm{X}, \mu)\) whose class \(\widetilde{h}\) belongs to \(\mathrm{E}_{\widetilde{m}_{g}}\) and is orthogonal to G. We then have

\[
0 = (\widetilde {h} \mid \widetilde {h} ^ {\prime}) _ {\widetilde {m} _ {g}} = \int_ {\mathrm{X}} h \overline {{h}} ^ {\prime} (1 + | g | ^ {2}) d \mu
\]

for all \(h' \in F\) of class \(h'\) Given the hypothesis on F, this means that the measure \(h(1 + |g|^{2}) \cdot \mu\) is zero, hence h is zero \(\mu\) -almost everywhere since h is moderate.

Let \(p\) a real number \(\geqslant 1\) Let \(h\) be an element of \(\mathcal{L}_{\mathrm{K}}^{\infty}(\mathrm{X},\mu)\) . The operator \(\widetilde{m}_h\) of multiplication by \(h\) is an endomorphism of \(\mathrm{L}_{\mathrm{K}}^{p}(\mathrm{X},\mu)\) Assume that the set Y of \(x\in X\) such that \(h(x)=0\) is locally \(\mu\) -negligible. The endomorphism \(\widetilde{m}_h\) is then injective (Lemma 7 of IV, p. 186). Denote \(h^{-1}\) the function on X equal to 0 on Y and to \(x \mapsto 1/h(x)\) on \(\mathrm{X} - \mathrm{Y}\). The reciprocal partial operator \(\widetilde{m}_{h}^{-1}\) is the operator of multiplication by \(h^{-1}\) of \(\mathrm{L}_{\mathrm{K}}^{p}(\mathrm{X},\mu)\) in \(\mathrm{L}_{\mathrm{K}}^{p}(\mathrm{X},\mu)\) , that is, \(\widetilde{m}_{h}^{-1} = \widetilde{m}_{h^{-1}}\) Indeed, the image of \(\widetilde{m}_{h}\) is the space of classes of functions \(g \in \mathcal{L}_{\mathrm{K}}^{p}(\mathrm{X},\mu)\) of the form \(g = hf\) for \(f \in \mathcal{L}_{\mathrm{K}}^{p}(\mathrm{X},\mu)\) This condition is equivalent to \(g(x)/h(x) = f(x)\) for all \(x \in \mathrm{X} - \mathrm{Y}\) and \(g(x) = 0\) if \(x \in \mathrm{Y}\) This implies that the domain of \(\widetilde{m}_{h}^{-1}\) in \(\mathrm{L}_{\mathrm{K}}^{p}(\mathrm{X},\mu)\) is the domain of \(\widetilde{m}_{h^{-1}}\) , and that the equality \(\widetilde{m}_{h}^{-1} = \widetilde{m}_{h^{-1}}\) is valid.

In what follows, we will sometimes simply denote by \(m_{h}\) the partial multiplication operator by h on \(\mathrm{L}_{\mathrm{K}}^{p_{1}}(\mathrm{X},\mu)\) in \(\mathrm{L}_{\mathrm{K}}^{p_{2}}(\mathrm{X},\mu)\) .

\begin{enumerate}
\def\labelenumi{\arabic{enumi})}
\setcounter{enumi}{1}
\tightlist
\item
  Let E be a Hilbert space over K and \(\mathrm{B} = (e_i)_{i \in \mathrm{I}}\) an orthonormal basis of E. Let \((\lambda_i)_{i \in \mathrm{I}}\) a family of elements of K. Let D be the vector subspace of E consisting of the elements \(x \in E\) such that the family \((\lambda_i \langle e_i | x \rangle)_{i \in \mathrm{I}}\) is square-summable in K. The space D is dense in E since it contains the vector \(e_i\) for all \(i \in I\) . The partial operator u with domain D given by
\end{enumerate}

\[
x \mapsto \sum_ {i \in I} \lambda_ {i} \langle e _ {i} | x \rangle e _ {i}
\]

is called a partial diagonal operator in the basis B, and \((\lambda_{i})_{i\in I}\) is called the family of eigenvalues of u.

The operator \(u\) is closed. Indeed, let \((x_{n}, u(x_{n}))_{n \in \mathbf{N}}\) a sequence of elements of the graph of \(u\) which converges in \(\mathrm{E} \times \mathrm{E}\) and let \((x, y)\) its limit. We then have \(\langle e_{i} | x_{n} \rangle \to \langle e_{i} | x \rangle\) for all \(i \in \mathrm{I}\) and

\[
\langle e _ {i} | u (x _ {n}) \rangle = \lambda_ {i} \langle e _ {i} | x _ {n} \rangle \rightarrow \langle e _ {i} | y \rangle
\]

for all \(i \in I\) . Consequently, \(\lambda_{i} \langle e_{i} \mid x \rangle = \langle e_{i} \mid y \rangle\) for all \(i \in I\) , which proves that \(x \in D\) and \(u(x) = y\) that is, \(u\) is closed.

This example is in fact a special case of the preceding one, applied to the topological space X = I equipped with the discrete topology and the counting measure \(\mu\) , since E is identified with the space \(\ell^{2}(\mathrm{I}) = \mathrm{L}^{2}(\mathrm{I}, \mu)\) by the mapping \(x \mapsto (\langle e_{i} \mid x \rangle)_{i \in \mathrm{I}}\) (EVT, V, p. 23, cor. 2) and u is then identified with the multiplication operator \(m_{\lambda}\) , where \(\lambda\) is the function \(i \mapsto \lambda_{i}\) .

\begin{enumerate}
\def\labelenumi{\arabic{enumi})}
\setcounter{enumi}{2}
\tightlist
\item
  The set \(N_{R} = N \cup \{\infty, \omega\}\) is equipped with the total order described in VAR, R2, p. 10, such that \(n < \infty < \omega\) for all \(n \in N\) Let \(r \in N_{R}\) . Let \(n \in N\) and U an open subset of \(R^{n}\) Let \(k \in N\) such that \(k \leqslant r\) Let \((n_{\alpha})_{|\alpha| \leqslant k}\) a family of elements of \(\mathcal{C}^{r}(U)\) , where the multi-indices considered belong to \(N^{n}\) . The family \((n_{\alpha})\) defines a scalar differential operator D of order \(\leqslant k\) on U (cf.VAR, R2, 14.1.6, 14.1.4).
\end{enumerate}

For every integer m such that \(k \leqslant m \leqslant r\) , the differential operator D defines a linear map from \(\mathrm{C}^{m}(\mathrm{U})\) in \(\mathrm{C}^{m-k}(\mathrm{U})\) which sends \(f \in \mathrm{C}^{m}(\mathrm{U})\) onto

\[
\mathrm{D} (f) = \sum_ {| \alpha | \leqslant k} n _ {\alpha} \partial^ {\alpha} (f).
\]

The same formula defines a continuous linear map from \(\mathcal{D}(\mathrm{U})\) in \(\mathcal{D}'(\mathrm{U})\) (def. 2 of IV, p. 208).

DEFINITION 5. --- Let E be a vector subspace of \(\mathcal{D}'(\mathrm{U})\) containing \(\mathcal{D}(\mathrm{U})\) A differential operator associated with D on E is any partial operator on E that is an extension of the partial operator with domain \(\mathcal{D}(\mathrm{U})\) defined by \(\varphi \mapsto \mathrm{D}(\varphi)\) .

Suppose for example that the coefficients \(n_{\alpha}\) are bounded functions on U. Let \(\mu\) Lebesgue measure on U. If p is an element of \([1,+\infty[\) one can then define a differential operator on \(\mathrm{L}^{p}(\mathrm{U})\) associated with D whose domain is the Sobolev space \(\mathrm{W}^{k,p}(\mathrm{U})\) (no. 14 of IV, p. 221), since in this case one has \(n_{\alpha}\partial^{\alpha}(f)\in\mathrm{L}^{p}(\mathrm{U})\) for all \(f\in\mathrm{W}^{k,p}(\mathrm{U})\) and every \(|\alpha|\leqslant k\) .

